\section{The unbinned OmniFold procedure}
\label{sec:method}

\subsection{Algorithm}
\label{sec:method-algo}
OmniFold~\cite{Andreassen:2019cjw,Andreassen:2021zzk} unfolds by iterated
reweighting with machine-learned likelihood ratios, operating on per-event
observables rather than on a fixed histogram binning. Each iteration has two
steps:
\begin{enumerate}[itemsep=2pt]
  \item \textbf{Step 1 (detector space).} A classifier is trained to separate
        the measured data from the reconstructed simulation (the latter
        weighted by the current push-weights). Its calibrated output $p$
        yields the likelihood-ratio reweight $w = p/(1-p)$ on every
        reconstructed simulated event. The weight is \emph{pulled back} to
        truth level through the event pairing: each simulated event is a
        (truth, reco) pair, so the reco-space weight attaches directly to the
        truth record. Truth events with no reconstructed partner (the
        ``misses'') receive a step-1 weight from a regressor trained on the
        reconstructed subsample, $w(\mathrm{truth\ features})$ --- the native
        OmniFold acceptance treatment, with no separate efficiency
        correction.
  \item \textbf{Step 2 (truth space).} A second classifier is trained to
        separate the original truth-level simulation from the same sample
        carrying the pulled-back step-1 weights, yielding a smooth truth-space
        reweight that becomes the next iteration's push-weight.
\end{enumerate}
After $n$ iterations the push-weights define an efficiency-corrected,
truth-level reweighted simulation ensemble; histogramming it (in any binning,
chosen after the fact) gives the unfolded spectrum. The regulariser is the
iteration count together with the model capacity of the learners.

\subsection{Estimators and configuration}
\label{sec:method-config}
We use gradient-boosted decision trees (GBDTs) rather than the neural networks
of the original OmniFold papers; \S\ref{sec:validation} validates this choice
with calibration tests and MLP/transformer cross-checks. The engine
(\texttt{unbinned\_unfolding/python/omnifold.py}) supports interchangeable
backends; the production configuration is \textbf{LightGBM} with
$100$ trees, $8$ leaves, learning rate $0.1$ (identical settings for the
step-1 classifier, step-2 classifier, and miss regressor). Two alternative
backends --- \texttt{exact} (sklearn \texttt{GradientBoosting}) and
\texttt{hist} (sklearn \texttt{HistGradientBoosting}, matched 8-leaf
configuration) --- are retained for method-dependence studies
(\S\ref{sec:results-tension}). The production iteration count is
$n=5$; doubling to $10$ moves the total cross section by
\SI{0.026}{\percent} (\S\ref{sec:validation}). The central value is a
single-run unfold with pinned seeds; an ensemble-mean central value (averaging
independent reweighting trials) agrees with it to \SI{0.28}{\percent}, and ML
stochasticity is budgeted explicitly via the seedscan band
(\S\ref{sec:systematics}).

\subsection{MINERvA driver and invariants}
\label{sec:method-driver}
The MINERvA driver is \texttt{2d-unfolding/unfold\_2d\_omnifold\_unbinned.py};
the higher-dimensional driver
(\texttt{nd-unfolding/unfold\_nd\_omnifold\_unbinned.py}) generalises it to an
axis list whose first two entries are always $(\pT,\pPar)$ (carrying the
fiducial gate and the per-$\pT$ flux) and \texttt{import}s the same 2D
helpers. The unfolded observables are the muon $(\pT,\pPar)$, extended by
$\Eavail$, $q_{3}$, $W$ in the higher-dimensional runs; the phase space is
that of \S\ref{sec:selection}. The pipeline enforces the following contract:
\begin{itemize}[itemsep=2pt]
  \item measured data are masked to the reconstructed phase space before
        step-1 training;
  \item the reco-space target is the POT-scaled, non-negative
        $\mathrm{data}-\mathrm{background}$, with signal \emph{fakes}
        (reconstructed in phase space, truth outside) included in the
        subtraction so both sides of the step-1 training are fake-free;
  \item the unfolded yield is filled with
        $w_{\mathrm{push}}\times w_{\mathrm{truth}}$ (the MnvTune
        central-value weight rides along), and is \emph{not} divided by an
        efficiency histogram --- OmniFold's miss treatment already returns an
        efficiency-corrected truth spectrum;
  \item closure runs restrict pseudo-data to events passing both reco and
        truth selections on the same weight footing as production.
\end{itemize}

\subsection{Cross-section extraction}
\label{sec:method-xsec}
Unfolded truth-level counts are converted to the differential cross section
via
\begin{equation*}
  \frac{d^{n}\sigma}{d x_{1}\cdots d x_{n}}
  = \frac{U}{c\,\Phi\,N\,\mathrm{POT}\,\Delta x_{1}\cdots\Delta x_{n}}
    \times 10^{4},
\end{equation*}
where $U$ is the unfolded yield, $c$ the (unit) completeness, $\Phi$ the
per-$\pT$ flux integral (broadcast over the other axes), $N$ the fiducial
nucleon count ($3.2353\times10^{30}$), and POT the data exposure
($1.057\times10^{21}$). The implementation is
\texttt{extract\_cross\_section\_2d} and its N-dimensional generalisation
\texttt{xsec\_nd.py} (self-tested to reproduce the frozen lower-dimensional
extractions to $<10^{-12}$). The flux normalisation enters as
$\sigma\propto1/\Phi$; correctly varying $\Phi$ per flux universe was the
dominant systematic fix (\S\ref{sec:systematics}).

\subsection{Binning}
\label{sec:method-binning}
The reported 2D binning (14 $\pT$ $\times$ 16 $\pPar$ bins; 205 of 224
reported) follows \texttt{minerva\_paper\_anc/bin\_mapping.txt}; the 19
unreported diagonal cells have
$\pT^{\mathrm{hi}}/\pPar^{\mathrm{lo}}>\tan20^{\circ}$. The extension axes use
$\Eavail$: $[0,0.1,0.2,0.4,0.8,1.5,3.0,100]\,\si{GeV}$;
$q_{3}$: $[0,0.2,0.4,0.6,0.8,1.2,2.0,100]\,\si{GeV}$;
$W$: $[0,1.1,1.4,1.8,2.2,3.0,100]\,\si{GeV}$ --- each ending in a wide catch
bin so the marginals capture the full CC-inclusive tails (a truncated top edge
would break the marginalisation anchors of \S\ref{sec:3d-anchor}). Because the
binning is applied only after unfolding, these choices can be revisited
without re-running the unfold. The full edge table and reported-bin mapping
are recorded in \texttt{2D\_OMNIFOLD\_REFERENCE.md} and
App.~\ref{app:statmethods}.

Although OmniFold uses no response matrix, the detector smearing the unfold
must undo is conveniently summarised by the truth$\to$reco migration on the
analysis binning (Fig.~\ref{fig:migration}): the diagonal purity is
$\sim0.6$ (median) on every axis, i.e.\ bin-to-bin migrations are
substantial and a bin-by-bin correction would be badly biased --- the
regime unfolding is for.

\begin{figure}[H]
  \centering
  \includegraphics[width=0.99\textwidth]{migration_resolution.png}
  \caption{Truth$\to$reco migration of the simulated signal on the analysis
  binning (top: column-normalised maps; bottom: diagonal purity per truth
  bin) for the five observables. Median diagonal purity is
  $0.58$--$0.61$ per axis. Shown for illustration only --- the unfold itself
  is unbinned.}
  \label{fig:migration}
\end{figure}
